\documentclass[11pt]{article}

\usepackage[T1]{fontenc}
\usepackage[utf8]{inputenc}
\usepackage{lmodern}
\usepackage[margin=1in]{geometry}
\usepackage{amsmath,amssymb}
\usepackage{enumitem}
\usepackage[hidelinks]{hyperref}

\setlist[enumerate]{leftmargin=*, itemsep=0.65em, topsep=0.65em}
\emergencystretch=2em

\title{Mathematical Revision Log}
\author{}
\date{}

\begin{document}

\maketitle

The manuscript has been revised in the following mathematically substantive ways. Purely typographical, layout, indexing, labeling, and hyperlink changes are omitted.

\begin{enumerate}
    \item Two mathematical gaps are fixed. The first is about Lemma~2.4.4. The original proof is a reproduction of Darvas--Di Nezza--Lu's argument in their survey paper. I realized that the proof contained a gap. The gap is fixed now using a transfinite induction. The second is in the proof of Theorem~8.4.1, the proof was reproduced from my joint paper with Darvas. The proof contains again a gap. The gap is fixed in the revised version, and we will update the arXiv version of our paper soon.
    \item Section~11.3 is removed for two reasons: First, the original proof contains a fixable gap, but the fix would take a lot of space. Second, the result is not used in the rest of the book.
    \item I carried out a systematic proofread of Guedj--Zeriahi's monograph \emph{Degenerate Complex Monge--Ampère Equations} with the help of AI. It turns out that the proofs of a lot of results are not complete/correct. I included several footnotes to explain these issues.
    \item With the help of Claude and ChatGPT, all figures are redrawn using TikZ.
    \item Two more examples (12.1.3 and 12.1.4) are included into the toric part to explain the pathologies more clearly.
    \item A lot of new references are added, including Rashkowskii's paper mentioned by the editor.
    \item At the end of Section~2.4, the proofs of several results related to the $P$-envelope are reproduced.
    \item The proofs of a lot of results are expanded, so that they are more accessible to students.
    \item The introductions to the book, to each part, to each chapter are all rewritten. I hope the concerns of the editors are addressed in the new introductions.
\end{enumerate}

There is one thing in R3 which I did not follow: I insist on calling the volume of currents $\mathrm{vol}$ instead of $\mathrm{vol}_{\mathcal{I}}$. This is mainly because the former notation is standard when denoting the volume of line bundles, and I want to keep the notation consistent. I hope this is not a problem.

I did not take R5 seriously, since similar to the previous report, the referee (presumably still José) is apparently not familiar with the subject, and the criticism is mostly based on personal feud instead of mathematical substance. 

\end{document}
